% ============================================================================
% CAPÍTULO 8 — Next.js\index{Next.js}
% Texto completo (rodada editorial 2)
% ============================================================================
\chapter{Next.js: O Framework Full Stack do React}
\label{cap:nextjs}

\begin{caixaconceito}
\textbf{Next.js} é o framework full stack mais usado sobre React
\cite{nextjs-docs}: ele adiciona roteamento por arquivos, renderização no
servidor (SSR/SSG\index{SSG}/ISR\index{ISR}), Server Components\index{Server Components}, Server Actions e otimizações
automáticas — a versão 16 \cite{nextjs16} trouxe Turbopack estável e Cache
Components. É o framework que o mercado pede nas vagas de full stack.
\end{caixaconceito}

Ao final deste capítulo, você será capaz de:
\begin{objetivos}
  \item Criar um projeto Next.js com o App Router\index{App Router} e entender sua estrutura;
  \item Construir páginas, layouts, rotas dinâmicas e estados de carregamento/erro;
  \item Distinguir Server Components de Client Components\index{Client Components} (\texttt{"use client"});
  \item Escolher entre SSR\index{SSR}, SSG e ISR com critério técnico;
  \item Navegar com \texttt{Link} e otimizar imagens e fontes;
  \item Entregar o projeto do capítulo: um blog com MDX\index{MDX} e ISR.
\end{objetivos}

\section{O que o Next.js resolve}

O React puro (capítulo~\ref{cap:react}) renderiza no navegador: o servidor
envia uma página quase vazia e o JavaScript monta tudo — o que é ruim para
SEO (buscadores veem pouco conteúdo) e para o primeiro acesso (o usuário
espera o bundle baixar para ver algo).

O Next.js resolve isso renderizando \textbf{HTML no servidor}: a página chega
pronta, o conteúdo é indexável e o primeiro acesso é rápido. Depois, o
JavaScript ``hidrata'' a página, devolvendo a interatividade do React. E mais:
com Server Components e Server Actions, o mesmo projeto é \emph{frontend e
backend ao mesmo tempo} — sem servidores separados para manter.

\begin{caixaconceito}
\textbf{Hidratação (hydration)}: o HTML renderizado no servidor é
``reconectado'' ao React no navegador — os event handlers são anexados aos
elementos já existentes, sem recriar a página. É a ponte entre ``HTML pronto''
e ``aplicação interativa''.
\end{caixaconceito}

\section{Criando o projeto e entendendo o App Router}

\begin{lstlisting}[style=livro, language=Bash,
  caption={Criando um projeto Next.js.}, label={list:nextjs-create}]
npx create-next-app@latest meu-blog
cd meu-blog
npm run dev        # desenvolvimento em http://localhost:3000
npm run build      # build de produção
npm start          # roda o build de produção
\end{lstlisting}

O comando faz perguntas (TypeScript, ESLint, Tailwind, App Router...).
Responda \emph{sim} para TypeScript e App Router — a stack deste livro.

\begin{lstlisting}[style=livro, language=Bash,
  caption={Estrutura de pastas do App Router.}, label={list:nextjs-estrutura}]
app/
  layout.tsx        # layout raiz (html, body, fontes)
  page.tsx          # rota "/"
  blog/
    page.tsx        # rota "/blog"
    [slug]/
      page.tsx      # rota dinâmica "/blog/meu-artigo"
      not-found.tsx # 404 da rota dinâmica
  api/
    route.ts        # API Route (backend dentro do app)
components/
  Header.tsx
lib/
  posts.ts          # funções de dados
\end{lstlisting}

\begin{caixaconceito}
No App Router, \emph{arquivos = rotas}: \texttt{page.tsx} define a página,
\texttt{layout.tsx} o layout compartilhado, \texttt{loading.tsx} o estado de
carregamento e \texttt{error.tsx} o tratamento de erros. Pasta, arquivo, rota
— uma convenção única para todo o app, sem biblioteca de roteamento separada.
\end{caixaconceito}

\section{Layouts e rotas dinâmicas}

Um \texttt{layout.tsx} envolve as páginas de sua pasta — e \emph{persiste}
entre navegações (o React mantém o estado do layout enquanto você troca de
página):

\begin{lstlisting}[style=livro, language=TypeScript,
  caption={Layout raiz com fontes e metadata.}, label={list:nextjs-layout}]
import type { Metadata } from "next";

export const metadata: Metadata = {
  title: { default: "Meu Blog", template: "%s | Meu Blog" },
  description: "Blog sobre desenvolvimento full stack",
};

export default function LayoutRaiz({ children }: { children: React.ReactNode }) {
  return (
    <html lang="pt-BR">
      <body>
        <header>{/* navbar */}</header>
        <main>{children}</main>
        <footer>{/* rodapé */}</footer>
      </body>
    </html>
  );
}
\end{lstlisting}

Rotas dinâmicas usam colchetes na pasta. A partir do Next.js 15/16,
\texttt{params} é uma \textbf{Promise} — você precisa \texttt{await}:

\begin{lstlisting}[style=livro, language=TypeScript,
  caption={Rota dinâmica com params assíncrono.}, label={list:nextjs-dinamica}]
interface PaginaParams {
  params: Promise<{ slug: string }>;
}

export default async function PaginaDoPost({ params }: PaginaParams) {
  const { slug } = await params; // params é Promise no Next 15/16
  const post = await obterPost(slug);

  if (!post) {
    return <h1>Post não encontrado</h1>;
  }

  return (
    <article>
      <h1>{post.titulo}</h1>
      <time dateTime={post.data}>{post.dataFormatada}</time>
      <p>{post.resumo}</p>
    </article>
  );
}
\end{lstlisting}

\section{Server Components vs.\ Client Components}

Por padrão, componentes no Next.js são \textbf{Server Components}: rodam no
servidor, podem acessar o banco diretamente e \emph{não enviam JavaScript ao
navegador}. Quando um componente precisa de interatividade (estado, eventos,
hooks), adicione \texttt{"use client"}:

\begin{lstlisting}[style=livro, language=TypeScript,
  caption={Server vs Client Components.}, label={list:nextjs-sc-cc}]
// Server Component (padrão) — pode buscar dados e acessar o banco
import { obterPosts } from "@/lib/posts";

export default async function PaginaInicial() {
  const posts = await obterPosts(); // direto no servidor!
  return <ListaDePosts posts={posts} />;
}

// Client Component — interativo (arquivo separado, com "use client")
"use client";
import { useState } from "react";

export function ContadorDeLikes({ postId }: { postId: string }) {
  const [likes, setLikes] = useState(0);
  return (
    <button onClick={() => setLikes((l) => l + 1)}>
      {likes} curtidas
    </button>
  );
}
\end{lstlisting}

\begin{caixadica}
Regra prática: \emph{server first}. Só desça para o cliente o que realmente
precisa de interatividade. Menos JavaScript no navegador = site mais rápido,
melhor SEO e custo menor. Essa mentalidade — ``o que pode rodar no servidor,
roda no servidor'' — é o que diferencia seniors em 2026.
\end{caixadica}

\begin{caixaarmadilha}
Um Server Component \emph{não pode} usar \texttt{useState} nem
\texttt{useEffect} — se você tentar, o compilador avisa. E um Client Component
\emph{não pode importar} um Server Component (mas pode recebê-lo como
\texttt{children}). A regra do ``arquivo com \texttt{"use client"}'' fica no
topo do arquivo, e tudo o que ele importa vira client também.
\end{caixaarmadilha}

\section{Renderização: SSR, SSG e ISR}

Três estratégias cobrem todos os cenários de atualização de conteúdo:

\begin{itemize}[leftmargin=*, itemsep=0.4em]
  \item \textbf{SSR} (Server-Side Rendering): HTML gerado a \emph{cada
  requisição} — para dados que mudam o tempo todo (painel com métricas ao vivo);
  \item \textbf{SSG} (Static Site Generation): HTML gerado \emph{no build} —
  para páginas quase imutáveis (página ``Sobre'', documentação);
  \item \textbf{ISR} (Incremental Static Regeneration): SSG com
  \emph{revalidação} — a página é estática, mas se regenera a cada N segundos
  ou sob demanda. É o padrão ideal para blogs e catálogos.
\end{itemize}

A Figura~\ref{fig:isr-fluxo} mostra o ciclo do ISR: enquanto o cache está
válido, a página estática é servida instantaneamente; quando expira, ela é
regenerada em segundo plano e o cache é substituído — o usuário nunca espera
pela geração.

\begin{figure}[htbp]
  \centering
  \diagramatikz{%
  \begin{tikzpicture}[
      box/.style={draw=primaria, fill=azulclaro, rounded corners=2pt,
        align=center, inner sep=4pt, font=\footnotesize},
      dec/.style={diamond, draw=primaria, fill=azulclaro, aspect=1.8,
        align=center, inner sep=1pt, font=\footnotesize},
      seta/.style={-{Stealth[length=2.2mm]}, thick, color=secundaria},
      font=\footnotesize
    ]
    \node[box] (build) at (0,0) {Build};
    \node[box] (static) at (3.0,0) {Página estática\\revalidate: 3600s};
    \node[dec] (d) at (6.4,0) {Cache\\válido?};
    \node[box] (cache) at (10.0,0) {Serve do\\cache};
    \node[box] (regen) at (6.4,-2.8) {Regenera em\\background};
    \node[box] (atual) at (10.0,-2.8) {Atualiza o\\cache};
    \draw[seta] (build) -- (static);
    \draw[seta] (static) -- (d);
    \draw[seta] (d) -- node[above] {sim} (cache);
    \draw[seta] (d) -- node[right] {não} (regen);
    \draw[seta] (regen) -- node[above] {nova página} (atual);
    \draw[seta] (atual.west) .. controls (9.0,-3.8) and (4.2,-3.8) ..
      (static.south) node[midway, below] {substitui};
  \end{tikzpicture}}
  \caption{Ciclo do ISR: estática no build, regenerada em background quando o
  cache expira.}
  \label{fig:isr-fluxo}
\end{figure}

\begin{lstlisting}[style=livro, language=TypeScript,
  caption={ISR: página estática que se atualiza.}, label={list:nextjs-isr}]
// Regenera a página a cada 1 hora (em segundo plano)
export const revalidate = 3600;

export default async function PaginaDoPost({ params }: {
  params: Promise<{ slug: string }>;
}) {
  const { slug } = await params;
  const post = await obterPost(slug);
  return <article>{/* conteúdo do post */}</article>;
}
\end{lstlisting}

\begin{caixaconceito}
No Next.js 16, o cache virou \emph{padrão por rota} (Cache Components) com
controle fino: \texttt{revalidate} para tempo, \texttt{revalidatePath} e
\texttt{revalidateTag} para invalidação sob demanda (você usará isso no
capítulo~\ref{cap:fullstack}). A pergunta que você deve sempre fazer: \emph{o
que muda? com que frequência? quem precisa ver na hora?}
\end{caixaconceito}

\section{Navegação, imagens e fontes otimizadas}

\begin{lstlisting}[style=livro, language=TypeScript,
  caption={Link, Image e fontes do Next.js.}, label={list:nextjs-otimizacao}]
import Link from "next/link";
import Image from "next/image";

export function Nav() {
  return (
    <nav>
      <Link href="/blog">Blog</Link>
      <Link href="/sobre">Sobre</Link>
    </nav>
  );
}

// Image otimiza formato (AVIF/WebP), tamanho e lazy-load automaticamente
<Image
  src="/capa-do-post.jpg"
  alt="Capa do artigo sobre Next.js"
  width={1200}
  height={630}
  priority
/>
\end{lstlisting}

\begin{caixadica}
\texttt{next/link} faz navegação \emph{client-side} (sem recarregar a página,
com preload da rota de destino) — por isso é preferível ao \texttt{<a>}
comum para rotas internas. \texttt{next/image} entrega as imagens no tamanho
certo para cada tela, no formato moderno e com \texttt{lazy loading} — os três
pilares da performance que você vai aprofundar no capítulo~\ref{cap:performance}.
\end{caixadica}

\section{SEO com Metadata}

Cada página pode declarar seus próprios metadados — dinamicamente ou não:

\begin{lstlisting}[style=livro, language=TypeScript,
  caption={Metadata dinâmica por página.}, label={list:nextjs-seo}]
import type { Metadata } from "next";export async function generateMetadata({ params }: { params: Promise<{ slug: string }>; }): Promise<Metadata> {
  const { slug } = await params;
  const post = await obterPost(slug);

  return {
    title: post.titulo,
    description: post.resumo,
    openGraph: {
      title: post.titulo,
      description: post.resumo,
      type: "article",
    },
  };
}
\end{lstlisting}

% ============================================================================
\section{Projeto do capítulo: blog com MDX e ISR}
\label{sec:projeto8}

\begin{caixaprojeto}
\textbf{Objetivo:} construir um blog completo onde cada artigo é um arquivo
MDX (Markdown + JSX), com listagem, página individual, ISR e tema visual.

\textbf{Critérios de aceite:}
\begin{itemize}[leftmargin=*, itemsep=0.2em]
  \item Lista de posts em \texttt{/blog} com título, resumo e data;
  \item Rota dinâmica \texttt{/blog/[slug]} com \texttt{generateStaticParams} e \texttt{revalidate};
  \item Conteúdo em MDX com pelo menos um componente JSX customizado dentro do artigo;
  \item \texttt{loading.tsx} e \texttt{not-found.tsx} para a rota dinâmica;
  \item Metadata dinâmica (\texttt{generateMetadata}) para SEO;
  \item Deploy na Vercel (grátis) com o projeto público.
\end{itemize}
\end{caixaprojeto}

Para gerar as páginas estáticas do blog, use \texttt{generateStaticParams} —
o Next.js pré-renderiza cada \texttt{slug} no build:

\begin{lstlisting}[style=livro, language=TypeScript,
  caption={generateStaticParams + not-found.}, label={list:nextjs-params}]
export async function generateStaticParams() {
  const posts = await obterTodosPosts();
  return posts.map((post) => ({ slug: post.slug }));
}

export default async function PaginaDoPost({ params }: {
  params: Promise<{ slug: string }>;
}) {
  const { slug } = await params;
  const post = await obterPost(slug);

  if (!post) {
    return <h1>404 — Post não encontrado</h1>;
  }
  return <article>{/* conteúdo */}</article>;
}
\end{lstlisting}

\begin{caixadica}
O blog será a vitrine de tudo o que você aprendeu até aqui — e a base do seu
portfólio. No capítulo~\ref{cap:fullstack}, ele ganha banco de dados,
comentários e autenticação.
\end{caixadica}

\section{Exercícios de fixação}

\begin{exercicios}
  \item Explique a diferença entre Server Components e Client Components e quando usar cada um.
  \item O que diferencia SSR, SSG e ISR? Dê um caso de uso real para cada.
  \item Como o Next.js transforma pastas em rotas? Cite três convenções de arquivo especiais.
  \item Escreva o comando para criar um projeto Next.js e a estrutura de pastas de uma rota dinâmica de produto.
  \item Por que \texttt{Link} é preferível a \texttt{<a>} para navegação interna?
  \item O que \texttt{generateStaticParams} faz e quando é necessário?
\end{exercicios}

\section{Desafios}

\begin{desafios}
  \item \textbf{Rotas dinâmicas}: crie a rota \texttt{/produtos/[id]} com \texttt{generateStaticParams}, \texttt{generateMetadata} e página de \emph{não encontrado}.
  \item \textbf{Layout aninhado}: crie um \texttt{layout.tsx} para a área \texttt{/dashboard} com sidebar, e uma rota de grupo \texttt{(marketing)} sem o layout do dashboard.
  \item \textbf{Server Action de contato}: implemente uma Server Action que recebe os dados do formulário de contato, valida e registra em log no servidor.
  \item \textbf{Notações de código}: crie um componente MDX customizado (ex.: \texttt{<Dica>}) e use-o em um artigo.
\end{desafios}

\section{Exercícios de nível sênior}

\begin{seniores}
  \item \textbf{Cache do Next.js 16}: explique as camadas de cache (Cache Components, rota, dados) e como \texttt{revalidatePath} e \texttt{revalidateTag} invalidam conteúdo sob demanda.
  \item \textbf{Streaming}: use \texttt{Suspense} para streaming de seções da página e explique o impacto em LCP (maior pintura de conteúdo).
  \item \textbf{Middleware e proxy:} implemente um \texttt{proxy.ts} que bloqueia requisições de países/IPs específicos ou aplica headers de segurança (base para o capítulo~\ref{cap:seguranca}).
  \item \textbf{Benchmark:} compare em Lighthouse a mesma página com \texttt{"use client"} pesado vs.\ Server Components e documente a diferença de JavaScript transferido.
\end{seniores}

\section{Armadilhas comuns}

\begin{itemize}[leftmargin=*, itemsep=0.4em]
  \item Marcar tudo com \texttt{"use client"} (perde todo o benefício do servidor);
  \item Buscar dados dentro de Client Components (deveria ser no servidor);
  \item Esquecer \texttt{await params} (Next.js 15/16 tornou params assíncrono);
  \item Usar \texttt{<img>} em vez de \texttt{next/image} (perde otimização);
  \item Ignorar \texttt{loading.tsx} — a primeira impressão de navegação fica ruim;
  \item Páginas 100\% estáticas com dados que mudam (dados defasados — use ISR).
\end{itemize}

\section{Resumo do capítulo}

\begin{itemize}[leftmargin=*, itemsep=0.3em]
  \item Next.js = React + roteamento + renderização no servidor + backend integrado;
  \item App Router: arquivos viram rotas; \texttt{layout/loading/error/not-found} são convenções;
  \item Server Components por padrão; \texttt{"use client"} só para interatividade;
  \item SSR, SSG e ISR cobrem todos os cenários de atualização de conteúdo;
  \item \texttt{next/link}, \texttt{next/image} e Metadata são otimizações embutidas;
  \item Projeto entregue: blog MDX com ISR publicado na Vercel.
\end{itemize}

\section{Referências oficiais para aprofundar}

\begin{itemize}[leftmargin=*, itemsep=0.3em]
  \item Documentação oficial do Next.js: \url{https://nextjs.org/docs}
  \item Anúncio do Next.js 16: \url{https://nextjs.org/blog/next-16}
  \item Curso oficial (App Router): \url{https://nextjs.org/learn}
  \item Documentação do MDX: \url{https://mdxjs.com}
  \item Guia de deploy na Vercel: \url{https://nextjs.org/learn-pages-router/basics/deploying-nextjs-app}
\end{itemize}
